\documentclass[12pt]{article}
\usepackage[T1]{fontenc}
\usepackage[utf8]{inputenc}
\usepackage[brazil]{babel}
\usepackage[margin=1in]{geometry}
\setlength{\parindent}{0pt}
\begin{document}
\section*{Tema 52}
Processo: PEDILEF 0505148-18.2010.4.05.8500/ SE\\
Situacao: Julgado\\
Ramo: DIREITO ADMINISTRATIVO\\
Questao: Saber se é devido restabelecimento do adicional de inatividade, suprimido dos proventos dos militares da reserva, por força da MP n. 2.131/2000.\\
Tese: A supressão do adicional de inatividade pela MP n. 2.131/2000 não afronta o princípio do direito adquirido, uma vez que não houve decesso remuneratório.\\
Fonte: https://www.cjf.jus.br/phpdoc/virtus/
\end{document}
